\documentclass[../main.tex]{subfiles}

\begin{document}

\section{Conclusions}\label{sec:conclusions}

A classical attitude control system for an aerodynamically unstable launch
vehicle at the max-$\bar q$ condition was designed and validated entirely in the
frequency domain.

\begin{itemize}
\item \textbf{Task 1.} A proportional--derivative attitude law with weak negative
drift feedback stabilises the rigid airframe (open-loop pole at
$+1.84$~rad/s) and meets the design targets $|GM|\approx 6$~dB,
$|PM|\approx 30^\circ$. The conditional-stability nature of the loop was
identified and accounted for in the interpretation of the margins.
\item \textbf{Task 2.} Introducing the TVC actuator, the transport delay and the
lightly damped bending mode destabilises the loop through a $+39$~dB resonance.
A gain-stabilising notch centred on $\omega_{BM}$ restores stability while
preserving the rigid margins, with negligible change to the nominal gust
response. Gain stabilisation was preferred to phase stabilisation because the
resonance is too large to be tamed by a phase shift alone.
\item \textbf{Task 3.} The fixed Task~2 controller is robust to $\pm 30\%$
independent variations of $\mu_\alpha$ and $\mu_c$: all four corners remain
stable. Aerodynamic uncertainty primarily erodes the rigid-body margin and
nominal performance, whereas control-effectiveness uncertainty primarily erodes
the bending margin.
\end{itemize}

The design illustrates the textbook trade-offs of launch-vehicle attitude
control~\cite{greensite1970,wie2008}: enough gain to stabilise the unstable
airframe, but not so much as to excite the structural modes, with the Nichols
chart providing a single, unified view of both constraints. A Simulink
block-diagram mirror of the closed loop reproduces the script results and is
documented in the accompanying \texttt{models/SIMULINK\_GUIDE.md}.

\end{document}
